\documentclass[pdflatex,sn-mathphys-num]{sn-jnl}

\usepackage{graphicx}
\usepackage{multirow}
\usepackage{amsmath,amssymb,amsfonts}
\usepackage{amsthm}
\usepackage{mathrsfs}
\usepackage[title]{appendix}
\usepackage{xcolor}
\usepackage{textcomp}
\usepackage{manyfoot}
\usepackage{booktabs}
\usepackage{algorithm}
\usepackage{algorithmicx}
\usepackage{algpseudocode}
\usepackage{listings}
\usepackage{url}

\theoremstyle{plain}
\newtheorem{proposicion}{Proposición}

\graphicspath{{figuras/}}

\raggedbottom

\begin{document}
\renewcommand{\abstractname}{Resumen}

\title[Selección dinámica para readmisión en diabéticos]{Selección dinámica de ensambles por competencia probabilística local para predecir readmisión temprana en diabéticos}

\author*[1]{\fnm{Joshua} \sur{Hincapié Llorente}}\email{hjoshua@uninorte.edu.co}

\author[1]{\fnm{Alejandro} \sur{Cantillo Escorcia}}\email{cantilloalejandro@uninorte.edu.co}

\affil*[1]{\orgdiv{Departamento de Matemáticas, Física y Ciencia de Datos}, \orgname{Universidad del Norte}, \orgaddress{\city{Barranquilla}, \country{Colombia}}}

\abstract{La readmisión a 30 días del alta es un indicador de calidad asistencial y un costo evitable, especialmente en pacientes diabéticos. La literatura que la predice con aprendizaje automático usa sobre todo clasificadores individuales y ensambles estáticos, que aplican la misma regla de combinación a todos los pacientes y alcanzan un desempeño modesto. La selección dinámica de ensambles, que pondera los clasificadores según la región en que se encuentra cada paciente, no ha sido explorada en este problema, y sus criterios de competencia basados en la exactitud local se degeneran cuando la clase positiva es minoritaria. Se propone DES-CPL, un ensamble dinámico que mide la competencia de cada modelo con la puntuación de Brier fuera de pliegue, en vecindarios estratificados por clase que excluyen al propio paciente, y los combina con pesos softmax de temperatura ajustable. Se demuestran una cota para su pérdida de Brier y la consistencia del estimador de competencia. Sobre Diabetes 130-US Hospitals (101,766 encuentros de 71,518 pacientes), DES-CPL alcanza un AUC-PR de [X.XXXX $\pm$ X.XXXX] frente a 0.2267 $\pm$ 0.0153 del mejor benchmark, XGBoost, con diferencias [significativas / no significativas] según [prueba 1, prueba 2 y prueba 3], y [resultado en los tres conjuntos adicionales]. El aporte es un criterio de competencia probabilístico, transferible a otros problemas de clasificación desbalanceada, evaluado con un protocolo que respeta la agrupación por paciente; [si supera: combinar localmente modelos que empatan en promedio mejora la identificación de pacientes en riesgo / si no: la combinación local no supera el techo de los modelos individuales].}

\keywords{Readmisión hospitalaria, Diabetes mellitus, Aprendizaje automático, Selección dinámica de ensambles, Desbalance de clases, Clasificación}

\maketitle

\section{Introducción}\label{sec:intro}

La readmisión hospitalaria no planificada dentro de los 30 días posteriores al alta es uno de los indicadores con los que los sistemas de salud miden la calidad de la atención. Jencks y colaboradores \cite{bib_jencks} estimaron que cerca de una quinta parte de los beneficiarios de Medicare dados de alta volvía al hospital dentro de ese plazo, con un costo anual cercano a 17 mil millones de dólares. En los pacientes diabéticos la cuestión es más aguda: la tasa de readmisión de quienes ingresan con ella es casi el doble de la de los pacientes sin diabetes \cite{bib_liu_2024}. Si al momento del alta fuera posible ordenar a los pacientes por su riesgo de reingreso, un servicio con capacidad limitada de seguimiento podría concentrar sus recursos en quienes más lo necesitan.

El problema que aborda este trabajo es que esa ordenación sigue siendo poco precisa. La revisión sistemática de Kansagara y colaboradores \cite{bib_kansagara} encontró que la mayoría de los modelos de readmisión obtenía un estadístico C entre 0.60 y 0.70. El conjunto Diabetes 130-US Hospitals, con diez años de registros de 130 hospitales estadounidenses \cite{bib_strack}, es la referencia habitual para estudiar el problema en pacientes diabéticos. Sobre él se han ensayado máquinas de vectores de soporte ajustadas con algoritmos genéticos y combinadas con SMOTE \cite{bib_cui_2018}, comparaciones entre bosques aleatorios, Naive Bayes y ensambles de árboles \cite{bib_shang_2021, bib_neto_2021}, baterías de diez o más algoritmos con selección de variables por metaheurísticas \cite{bib_liu_2024}, ensambles por apilamiento \cite{bib_lu_2022, bib_salim_2026} y, más recientemente, arquitecturas neuronales y Transformers \cite{bib_ram_2023, bib_garcia_mosquera_2025}. Con evaluación rigurosa el techo se mantiene: con validación cruzada anidada, un ensamble por apilamiento alcanza un AUC-ROC de 0.664 \cite{bib_salim_2026}, y las exactitudes cercanas al 90\% \cite{bib_neto_2021} apenas superan las de predecir siempre que el paciente no reingresa. Además, rara vez se valida respetando la estructura del conjunto, en el que un mismo paciente aporta varios encuentros. La pregunta que guía este trabajo es si una combinación que adapte el peso de cada modelo al perfil de cada paciente puede superar ese techo cuando la evaluación respeta la agrupación por paciente.

Nuestro benchmark previo confirma ese diagnóstico. De 108 configuraciones de siete modelos base, estimadas con validación cruzada anidada agrupada por paciente, la mejor, XGBoost \cite{bib_chen} con ADASYN, alcanzó un AUC-PR de 0.2267 $\pm$ 0.0153, y la prueba de Nemenyi \cite{bib_demsar} no la separó de Random Forest. Que modelos tan distintos empaten en promedio sugiere que cada uno acierta en regiones diferentes del espacio de pacientes, pero los ensambles usados hasta ahora en este problema son estáticos: aplican la misma regla de combinación a todos los pacientes y no aprovechan esa complementariedad local.

La selección dinámica de ensambles responde precisamente a esa limitación. Estima para cada consulta la competencia de cada clasificador en su vecindario dentro de un conjunto de validación, y elige o pondera a los más competentes \cite{bib_woods_1997, bib_ko_2008, bib_cruz_2018}. Existen variantes que podan los clasificadores no informativos \cite{bib_cruz_2019} y adaptaciones al desbalance que remuestrean el entrenamiento o ponderan las instancias del vecindario \cite{bib_garcia_2018, bib_hou_2020}. La propuesta más cercana a la nuestra es el modelo probabilístico de competencia de Woloszynski y Kurzynski \cite{bib_woloszynski_2011}, que estima con un clasificador de referencia aleatorizado la probabilidad de que cada clasificador acierte la etiqueta. Todas estas variantes comparten una limitación crítica aquí: miden el acierto de etiquetas y no la calidad de la probabilidad, que es lo que importa al ordenar pacientes por riesgo con el AUC-PR \cite{bib_saito}. Cuando solo el 11\% de los encuentros son positivos, muchos vecindarios no contienen ningún caso positivo y en ellos todos los clasificadores parecen igualmente competentes por predecir siempre la clase mayoritaria \cite{bib_cruz_2019}; además, ninguna contempla la agrupación de encuentros por paciente. En nuestra búsqueda en Web of Science y Scopus (257 y 83 registros sobre readmisión en pacientes diabéticos) no encontramos ningún estudio que aplicara selección dinámica de ensambles a este problema.

El objetivo de este trabajo es diseñar, fundamentar teóricamente y evaluar un ensamble dinámico cuyo criterio de competencia funcione con clases desbalanceadas, porque identificar mejor a los pacientes de mayor riesgo se traduce directamente en un mejor uso de los recursos de seguimiento. Para ello se propone DES-CPL, una selección dinámica de ensambles por competencia probabilística local, con cuatro contribuciones. Primero, la competencia de cada modelo base se mide con la puntuación de Brier \cite{bib_brier}, una regla de puntuación estrictamente propia \cite{bib_gneiting_2007}, calculada sobre predicciones fuera de pliegue, de modo que premia la calidad de la probabilidad y no solo el acierto de la etiqueta. Segundo, el vecindario se construye estratificado por clase, con un número fijo de vecinos positivos y negativos, y excluye los encuentros del mismo paciente, lo que garantiza que la competencia se mida también sobre la clase minoritaria y evita fugas de información. Tercero, los modelos se combinan con pesos softmax cuya temperatura recorre de forma continua el espectro entre el promedio estático y la selección del mejor modelo local, y se ajusta dentro de la validación anidada. Cuarto, se demuestra que la pérdida de Brier del ensamble en cada punto está acotada por la combinación convexa de las pérdidas de sus miembros, y que el estimador de competencia es consistente bajo las condiciones clásicas de la regresión por vecinos más cercanos \cite{bib_stone_1977}.

DES-CPL se implementa desde cero, compatible con scikit-learn, sobre las mejores configuraciones del benchmark previo, y se evalúa con exactamente el mismo protocolo que ellas: las mismas particiones externas agrupadas por paciente, la misma semilla y el mismo presupuesto de optimización bayesiana. La superioridad se contrasta con tres pruebas de hipótesis fijadas antes de los experimentos, entre ellas la de DeLong \cite{bib_demsar, bib_delong}, se replica en tres conjuntos adicionales de clasificación desbalanceada y se complementa con un estudio de ablación.

El resto del artículo se organiza así. La Sección~2 revisa el trabajo relacionado; la Sección~\ref{sec:metodologia} describe los datos, la formulación de DES-CPL, sus propiedades teóricas y el protocolo experimental; la Sección~4 presenta los resultados; la Sección~5 los discute frente a la literatura y expone las limitaciones, y la Sección~6 resume las conclusiones.

\section{Materiales y métodos}\label{sec:metodologia}

\subsection{Formulación del problema y notación}\label{subsec:formulacion}

Sea $\mathcal{D}=\{(x_i, y_i, g_i)\}_{i=1}^{n}$ la muestra de entrenamiento, donde $x_i \in \mathcal{X} \subseteq \mathbb{R}^{p}$, con $p=152$, describe el encuentro hospitalario $i$ después del preprocesamiento, $y_i \in \mathcal{Y}=\{0,1\}$ indica si el paciente reingresó antes de 30 días y $g_i$ identifica al paciente. Un mismo paciente puede aportar varios encuentros, de modo que las observaciones son independientes entre pacientes pero no dentro de ellos. La tarea es aprender una función $f:\mathcal{X}\to[0,1]$ que estime la probabilidad de reingreso $\eta(x)=P(Y=1\mid X=x)$ y que ordene correctamente a pacientes que no se observaron durante el entrenamiento.

Se dispone de un conjunto de $M$ clasificadores base $\{f_1,\dots,f_M\}$, cada uno con salida probabilística $f_m(x)\in[0,1]$. La pérdida con la que se mide la calidad de una probabilidad $q$ frente a la etiqueta $y$ es la puntuación de Brier \cite{bib_brier},
\begin{equation}
\ell(q, y) = (q - y)^2,
\label{eq:brier}
\end{equation}
una regla de puntuación estrictamente propia: su esperanza se minimiza solo cuando $q$ coincide con la probabilidad verdadera \cite{bib_gneiting_2007}. Por la Ecuación~\eqref{eq:brier}, un clasificador obtiene poca pérdida solo si asigna probabilidades altas a los positivos y bajas a los negativos, y no basta con que acierte la etiqueta. La métrica de evaluación es distinta de esta pérdida interna y se describe en la Sección~\ref{subsec:evaluacion}. La Tabla~\ref{tab:notacion} resume la notación.

\begin{table}[h]
\caption{Notación utilizada en la formulación de DES-CPL.}\label{tab:notacion}
\begin{tabular}{@{}ll@{}}
\toprule
Símbolo & Significado \\
\midrule
$x_i,\ y_i,\ g_i$ & Predictores, etiqueta y paciente del encuentro $i$ \\
$f_m$ & Clasificador base $m$, con $m=1,\dots,M$ \\
$\tilde{f}_m(x_i)$ & Predicción fuera de pliegue de $f_m$ para el encuentro $i$ \\
$\phi(x)$ & Perfil de salidas: vector de las $M$ predicciones para $x$ \\
$\mathcal{N}^{+}_{k}(x),\ \mathcal{N}^{-}_{k}(x)$ & $k$ vecinos positivos y $k$ negativos de $x$ \\
$C_m(x)$ & Competencia local de $f_m$ en $x$ (Brier balanceado) \\
$\tau$ & Temperatura de la ponderación softmax \\
$w_m(x)$ & Peso de $f_m$ en la combinación para $x$ \\
$\hat{f}(x)$ & Probabilidad final del ensamble \\
\botrule
\end{tabular}
\end{table}

\subsection{Datos y preprocesamiento}\label{subsec:datos}

\textbf{Conjunto principal.} Se usa \emph{Diabetes 130-US Hospitals for Years 1999--2008} \cite{bib_strack}, publicado en el repositorio UCI con licencia CC BY 4.0 \cite{bib_uci}. Contiene 101{,}766 encuentros de 71{,}518 pacientes en 130 hospitales estadounidenses, descritos por 50 atributos demográficos, administrativos, diagnósticos CIE-9, de laboratorio y de 23 medicamentos antidiabéticos. La variable \texttt{readmitted} se dicotomiza en reingreso antes de 30 días frente al resto. Se excluyen 2{,}426 encuentros con fallecimiento o cuidados paliativos al alta, en los que no puede observarse un reingreso, y quedan 99{,}340 encuentros. La clase positiva representa el 11.39\%, una razón de desbalance de 1 a 7.8.

\textbf{Preprocesamiento.} Se descartan \texttt{weight}, con un 98.5\% de faltantes, y nueve fármacos casi constantes. La ausencia de \texttt{race} (2.25\%) se conserva como categoría propia, los niveles poco frecuentes se agrupan y los diagnósticos CIE-9 se consolidan en grupos clínicos. Las variables numéricas se escalan con mediana e intervalo intercuartílico, porque todas rechazan la normalidad y tres tienen más del 50\% de ceros. Las categóricas se codifican \emph{one-hot}, lo que deja 152 columnas.

\textbf{Particiones y control de fuga.} El conjunto se divide por paciente en 79{,}473 encuentros de entrenamiento y 19{,}867 de prueba, con intersección de pacientes vacía y prevalencias de 11.389\% y 11.391\%. La partición es por paciente porque los encuentros de un mismo individuo comparten casi por completo sus predictores ($\eta^2=0.99$ en la edad): particionar por filas inflaba el AUC-PR de Random Forest en 0.0127, una cifra del mismo orden que las diferencias que se quieren detectar. Todo ajuste con parámetros aprendidos (imputación, escalado, codificación, remuestreo y las predicciones fuera de pliegue de DES-CPL) se calcula únicamente con los datos de entrenamiento de cada pliegue.

\textbf{Conjuntos adicionales.} Para evaluar si el mecanismo se transfiere a otros problemas de clasificación binaria desbalanceada, DES-CPL y los modelos de referencia se evalúan además en tres conjuntos públicos del repositorio UCI, de naturaleza distinta: \emph{CDC Diabetes Health Indicators} (253{,}680 encuestas, 13.9\% de positivos), \emph{Bank Marketing} (45{,}211 contactos, 11.7\% de positivos) \cite{bib_moro_2014} y \emph{Default of Credit Card Clients} (30{,}000 clientes, 22.1\% de positivos) \cite{bib_yeh_2009}. Ninguno tiene identificador de individuo, de modo que en ellos la partición es estratificada simple; el preprocesamiento es análogo (escalado robusto y codificación \emph{one-hot} ajustados dentro de cada pliegue).

\subsection{Modelo propuesto: DES-CPL}\label{subsec:propuesto}

\textbf{Idea central.} DES-CPL decide, para cada paciente nuevo, cuánto confiar en cada clasificador base según qué tan buenas fueron sus probabilidades en pacientes parecidos. La Figura~\ref{fig:descpl} resume sus cuatro componentes: un conjunto de selección dinámica construido con predicciones fuera de pliegue, una región de competencia estratificada por clase y que excluye al propio paciente, una medida de competencia basada en la puntuación de Brier y una combinación con pesos softmax. El marco de selección dinámica con vecinos más cercanos se reutiliza de la literatura \cite{bib_ko_2008, bib_cruz_2018}; la medida de competencia, la construcción del vecindario y la regla de ponderación son nuevas.

\begin{figure}[t]
\centering
\includegraphics[width=\linewidth]{fig_descpl_diagrama.pdf}
\caption{Esquema de DES-CPL. Las predicciones fuera de pliegue de los $M$ modelos base forman el conjunto de selección dinámica. Para cada consulta se buscan $k$ vecinos positivos y $k$ negativos en el espacio de perfiles de salida, excluyendo los encuentros del mismo paciente; con ellos se calcula la puntuación de Brier balanceada de cada modelo, que se convierte en pesos mediante una softmax con temperatura $\tau$.}
\label{fig:descpl}
\end{figure}

\textbf{Componente 1: conjunto de selección con predicciones fuera de pliegue.} La competencia de un modelo no puede medirse sobre los mismos encuentros con los que se entrenó, porque ahí sus probabilidades son optimistas y todos los modelos parecerían competentes. Por eso cada $f_m$ se reentrena en una validación cruzada agrupada por paciente sobre $\mathcal{D}$, y para cada encuentro se guarda la predicción del modelo que no lo vio, $\tilde{f}_m(x_i)$. El conjunto de selección es $\mathcal{S}=\{(\phi(x_i), y_i, g_i)\}_{i=1}^{n}$, con
\begin{equation}
\phi(x) = \big(\operatorname{logit} \bar{f}_1(x), \dots, \operatorname{logit} \bar{f}_M(x)\big) \in \mathbb{R}^{M},
\label{eq:perfil}
\end{equation}
donde, para los encuentros de $\mathcal{S}$, $\bar{f}_m(x_i)=\tilde{f}_m(x_i)$, y para una consulta nueva $\bar{f}_m(x)$ es el promedio de los $K$ modelos de pliegue, de modo que ambos perfiles provienen de modelos entrenados con el mismo tamaño de muestra y son comparables. Las probabilidades se recortan a $[10^{-6},\,1-10^{-6}]$ antes del logit, porque modelos como el árbol o Naive Bayes producen ceros y unos exactos. La distancia se mide en este espacio de perfiles de salida y no en el de los 152 predictores, que es disperso y de alta dimensión. Dos pacientes con perfiles cercanos son pacientes que los modelos consideran parecidos, que es justamente la noción de vecindad relevante para estimar competencia \cite{bib_cruz_2018}. La transformación logit separa las probabilidades pequeñas, que es donde se concentra la clase positiva.

\textbf{Componente 2: región de competencia estratificada y agrupada.} Con un 11\% de positivos, los $k$ vecinos más cercanos de una consulta contienen con frecuencia uno o ningún positivo, y en ese caso la competencia solo premia a los modelos que predicen probabilidades bajas. Para evitarlo, la región de competencia reúne los $k$ vecinos positivos más cercanos, $\mathcal{N}^{+}_{k}(x)$, y los $k$ negativos más cercanos, $\mathcal{N}^{-}_{k}(x)$, en distancia euclidiana sobre $\phi$. Además, cuando la consulta proviene del propio conjunto de selección, como ocurre al ajustar los hiperparámetros (Sección~\ref{subsec:entrenamiento}), se excluyen los encuentros con $g_i=g(x)$; de lo contrario, el paciente se usaría para evaluarse a sí mismo.

\textbf{Componente 3: competencia por Brier balanceado.} La competencia local del modelo $m$ es su puntuación de Brier promedio en la región, con igual peso para cada clase:
\begin{equation}
C_m(x) = \frac{1}{2k}\sum_{i\in\mathcal{N}^{+}_{k}(x)} \big(\tilde{f}_m(x_i) - 1\big)^2 + \frac{1}{2k}\sum_{i\in\mathcal{N}^{-}_{k}(x)} \tilde{f}_m(x_i)^2 .
\label{eq:competencia}
\end{equation}
Valores bajos indican mayor competencia. A diferencia de la exactitud local \cite{bib_ko_2008} o de la probabilidad de acierto \cite{bib_woloszynski_2011}, la Ecuación~\eqref{eq:competencia} distingue entre un modelo que asigna 0.45 a un positivo y otro que le asigna 0.05, aunque ambos fallen la etiqueta con el umbral de 0.5.

\textbf{Componente 4: ponderación softmax con temperatura.} Los pesos se obtienen como
\begin{equation}
w_m(x) = \frac{\exp\!\big(-C_m(x)/\tau\big)}{\sum_{j=1}^{M}\exp\!\big(-C_j(x)/\tau\big)}, \qquad
\hat{f}(x) = \sum_{m=1}^{M} w_m(x)\, f_m(x).
\label{eq:pesos}
\end{equation}
La temperatura $\tau>0$ recorre un espectro continuo: cuando $\tau\to\infty$ los pesos se igualan y DES-CPL coincide con el promedio estático, y cuando $\tau\to 0$ toda la masa se concentra en el modelo localmente más competente, que es la selección dinámica de un solo clasificador. En lugar de fijar uno de los dos extremos, $\tau$ se ajusta en los datos.

\textbf{Propiedades teóricas.} Las dos proposiciones siguientes justifican el diseño; las demostraciones completas se incluyen en el Apéndice~\ref{app:demostraciones}.

\begin{proposicion}[Descomposición local de la pérdida]\label{prop:cota}
Para todo $x$, toda etiqueta $y\in\{0,1\}$ y cualesquiera pesos $w_m(x)\geq 0$ con $\sum_m w_m(x)=1$,
\begin{equation}
\ell\big(\hat{f}(x), y\big) = \sum_{m=1}^{M} w_m(x)\,\ell\big(f_m(x), y\big) - \sum_{m=1}^{M} w_m(x)\big(f_m(x)-\hat{f}(x)\big)^2 .
\label{eq:ambiguedad}
\end{equation}
En consecuencia, la pérdida del ensamble nunca supera el promedio ponderado de las pérdidas de sus miembros, y la diferencia es la dispersión ponderada de sus predicciones.
\end{proposicion}

La Ecuación~\eqref{eq:ambiguedad} es la descomposición de ambigüedad de Krogh y Vedelsby \cite{bib_krogh_1995}, aplicada aquí con pesos que dependen de $x$. Su consecuencia práctica es que, si los pesos favorecen a los modelos con menor pérdida local, la pérdida del ensamble queda acotada por una cantidad que ya es pequeña.

\begin{proposicion}[Consistencia del estimador de competencia]\label{prop:consistencia}
Suponga que los pacientes son independientes e idénticamente distribuidos, que $\phi$ tiene distribución con densidad dentro de cada clase y que $E[(\tilde{f}_m(X)-Y)^2 \mid \phi(X)=u,\, Y=c]$ es continua en $u$. Si $k\to\infty$ y $k/n_c\to 0$ para cada clase $c$, donde $n_c$ es el número de encuentros de esa clase, entonces $C_m(x)$ converge en probabilidad al riesgo de Brier balanceado del modelo $m$ en la vecindad de $\phi(x)$.
\end{proposicion}

La Proposición~\ref{prop:consistencia} se sigue de los resultados de consistencia de la regresión por $k$ vecinos más cercanos \cite{bib_stone_1977}, aplicados por separado a cada clase. El supuesto de independencia se cumple entre pacientes, no entre encuentros; por eso la exclusión del propio paciente en el Componente~2 es necesaria para que el resultado aplique.

El Algoritmo~\ref{alg:descpl} resume el procedimiento.

\begin{algorithm}
\caption{DES-CPL: entrenamiento y predicción}\label{alg:descpl}
\begin{algorithmic}[1]
\Require Datos $\mathcal{D}=\{(x_i,y_i,g_i)\}$, modelos base $f_1,\dots,f_M$ con hiperparámetros fijos, vecinos por clase $k$, temperatura $\tau$, pliegues $K$
\Ensure Probabilidad $\hat{f}(x)$ para una consulta $x$
\Statex \textit{Entrenamiento}
\For{cada pliegue $(\mathcal{T}_r, \mathcal{V}_r)$ de \textsc{StratifiedGroupKFold}$(\mathcal{D}, K)$}
    \State Ajustar cada $f_m$ sobre $\mathcal{T}_r$ y guardar $\tilde{f}_m(x_i)$ para $i\in\mathcal{V}_r$
\EndFor
\State Construir $\mathcal{S}$ con los perfiles $\phi(x_i)$ de la Ec.~\eqref{eq:perfil}
\State Construir un índice de vecinos para los positivos de $\mathcal{S}$ y otro para los negativos
\State Reajustar cada $f_m$ sobre $\mathcal{D}$ completo
\Statex \textit{Predicción}
\State Calcular $f_m(x)$ con los modelos reajustados y el perfil $\phi(x)$ con el promedio de los $K$ modelos de pliegue
\State Obtener $\mathcal{N}^{+}_{k}(x)$ y $\mathcal{N}^{-}_{k}(x)$, excluyendo encuentros del paciente $g(x)$ si está en $\mathcal{S}$
\State Calcular $C_m(x)$ con la Ec.~\eqref{eq:competencia} y $w_m(x)$ con la Ec.~\eqref{eq:pesos}
\State \Return $\hat{f}(x)=\sum_m w_m(x)\,f_m(x)$
\end{algorithmic}
\end{algorithm}

\textbf{Complejidad.} Sea $T_m(n)$ el costo de ajustar el modelo $m$ con $n$ encuentros. El entrenamiento cuesta $O\big((K+1)\sum_m T_m(n)\big)$ por los $K$ ajustes fuera de pliegue y el ajuste final, más $O(nM\log n)$ por la construcción de los dos índices de vecinos. La predicción de una consulta cuesta la inferencia de los $M$ modelos más $O(M\log n)$ en promedio para la búsqueda de vecinos y $O(kM)$ para las competencias. La memoria adicional es $O(nM)$. DES-CPL no tiene parámetros entrenables propios: solo los hiperparámetros $k$ y $\tau$, más los de los modelos base.

\subsection{Entrenamiento y optimización}\label{subsec:entrenamiento}

El conjunto base está formado por las siete mejores configuraciones del benchmark previo, una por modelo: regresión logística, Naive Bayes, $k$ vecinos más cercanos, árbol de decisión, Random Forest, XGBoost y SVM lineal. Cada una conserva la estrategia de balanceo y los hiperparámetros que eligió su propio bucle interno en el mismo pliegue externo, de modo que DES-CPL no recibe información adicional. Las predicciones fuera de pliegue se calculan con $K=3$ pliegues agrupados por paciente.

Los hiperparámetros de DES-CPL son el número de vecinos por clase, $k\in\{5,\dots,50\}$, y la temperatura, $\tau\in[0.001, 0.3]$ en escala logarítmica, que se expresa en unidades de la puntuación de Brier: con diferencias típicas de competencia del orden de 0.01, este rango cubre desde una selección casi dura hasta un promedio casi uniforme. Se exploran con optimización bayesiana con estimador de Parzen estructurado en árbol \cite{bib_bergstra_tpe}, con 12 evaluaciones, el mismo presupuesto máximo que recibieron los modelos de referencia. Como DES-CPL ya dispone de predicciones fuera de pliegue, cada configuración se evalúa sin reentrenar los modelos base: cada encuentro de $\mathcal{S}$ se predice con vecinos que excluyen a su propio paciente, y se maximiza el AUC-PR resultante. Esta evaluación ocurre íntegramente dentro del pliegue externo de entrenamiento. La semilla es 42, la misma del benchmark previo, y se propaga a \texttt{numpy}, \texttt{scikit-learn}, XGBoost y Optuna. Los experimentos se ejecutan en CPU, con cuatro núcleos físicos.

\subsection{Modelos de referencia}\label{subsec:referencia}

DES-CPL se compara con tres grupos de modelos, todos entrenados con las mismas particiones, el mismo preprocesamiento y el mismo presupuesto de ajuste. El primero son los siete modelos base individuales, descritos y ajustados en el benchmark previo. El segundo son dos ensambles estáticos sobre el mismo conjunto base: el promedio simple de probabilidades y el apilamiento con un meta-clasificador logístico entrenado sobre las mismas predicciones fuera de pliegue. Estos ensambles permiten separar lo que aporta combinar modelos de lo que aporta combinarlos de forma local. El tercero son métodos de selección dinámica del estado del arte: KNORA-E y KNORA-U \cite{bib_ko_2008}, META-DES \cite{bib_cruz_2015} y RRC, el método probabilístico de Woloszynski y Kurzynski \cite{bib_woloszynski_2011}. Se usan sus implementaciones de referencia en DESlib \cite{bib_deslib}, con el mismo conjunto base, el mismo conjunto de selección y su número de vecinos ajustado con el mismo presupuesto de 12 evaluaciones.

\subsection{Protocolo de evaluación}\label{subsec:evaluacion}

\textbf{Validación.} Todos los modelos se evalúan con validación cruzada anidada de cinco pliegues externos y tres internos, estratificada y agrupada por paciente, con la semilla 42. Las particiones externas son exactamente las del benchmark previo, de modo que los resultados de los siete modelos base se reutilizan sin volver a ejecutarlos y cada modelo aporta cinco estimaciones emparejadas. La métrica reportada proviene siempre del bucle externo. El conjunto de prueba reservado se usa una sola vez, al final.

\textbf{Métricas.} La métrica principal es el AUC-PR, calculada como precisión media, porque con un 11\% de positivos refleja directamente la proporción de pacientes señalados que realmente reingresan y tiene como referencia aleatoria la prevalencia \cite{bib_saito}. Como métricas secundarias se reportan el AUC-ROC, para comparar con la literatura; la puntuación de Brier y el error de calibración esperado, para la calidad probabilística; y la exactitud balanceada y el coeficiente de correlación de Matthews con el umbral que maximiza $F_1$ en el bucle interno. Se reporta la media y la desviación estándar de cada métrica.

\textbf{Pruebas de hipótesis.} Antes de ejecutar los experimentos definitivos se fija la hipótesis nula de que DES-CPL no discrimina mejor que el mejor modelo de referencia, definido como el de mayor AUC-PR medio en el bucle externo. Se contrasta con tres pruebas unilaterales con $\alpha=0.05$, dos sobre el AUC-PR y una sobre el AUC-ROC:
\begin{enumerate}
\item la prueba $t$ corregida para validación cruzada repetida de Nadeau y Bengio \cite{bib_nadeau_2003}, sobre las cinco diferencias emparejadas de AUC-PR del bucle externo, con cuatro grados de libertad y la razón entre tamaños de validación y entrenamiento de $1/4$, porque los pliegues comparten datos de entrenamiento y una prueba $t$ ordinaria subestimaría la varianza. Se prefiere a la prueba de Wilcoxon porque, con cinco pares, esta no puede producir un valor $p$ menor que 0.0625;
\item la prueba de DeLong \cite{bib_delong} para el AUC-ROC en el conjunto de prueba, cuyos 19{,}867 encuentros compensan el poco poder de una prueba con cinco pliegues;
\item un intervalo bootstrap BCa de la diferencia de AUC-PR en el conjunto de prueba, remuestreando pacientes y no encuentros, con 2{,}000 réplicas.
\end{enumerate}
Las comparaciones contra los distintos modelos de referencia se corrigen por multiplicidad con el método de Holm, y se reporta el delta de Cliff como tamaño de efecto \cite{bib_cliff}. Entre los cuatro conjuntos de datos se aplican las pruebas de Friedman y Nemenyi \cite{bib_demsar}.

\textbf{Robustez.} Se mide la sensibilidad de DES-CPL a $k$ y $\tau$ fuera del óptimo; su estabilidad entre semillas, repitiendo solo DES-CPL y el mejor modelo de referencia con las semillas 7 y 2026, como se hizo en el benchmark previo; y su comportamiento cuando se invierte el 5\% y el 10\% de las etiquetas de entrenamiento, cuando el entrenamiento se reduce al 25\% y al 50\% de los pacientes y cuando se submuestrean los positivos hasta una prevalencia del 5\%.

\subsection{Diseño del estudio de ablación}\label{subsec:ablacion}

Cada variante cambia un solo componente y conserva todo lo demás, con el mismo protocolo, las mismas particiones y la prueba de Nadeau y Bengio:
\begin{enumerate}
\item \textbf{Competencia por exactitud local} en lugar de Brier. Pone a prueba si medir la calidad de la probabilidad, y no el acierto, aporta desempeño.
\item \textbf{Vecindario no estratificado}, con $2k$ vecinos sin distinguir clase. Pone a prueba si garantizar la presencia de positivos en la región mejora la estimación de competencia.
\item \textbf{Sin exclusión del paciente} al ajustar los hiperparámetros. Pone a prueba si la fuga intra-paciente sesga la selección de $k$ y $\tau$.
\item \textbf{Espacio de predictores} en lugar del espacio de perfiles de salida. Pone a prueba la elección de la Ecuación~\eqref{eq:perfil}.
\item \textbf{Extremos de la temperatura}: selección dura ($\tau\to 0$) y promedio estático ($\tau\to\infty$). Pone a prueba si la ponderación intermedia supera a ambos.
\end{enumerate}

\subsection{Disponibilidad de código y datos}\label{subsec:disponibilidad}

El código de DES-CPL, implementado desde cero como un clasificador compatible con \texttt{scikit-learn}, sus pruebas unitarias, los cuadernos del experimento y las versiones exactas de las dependencias se publican en \url{https://github.com/joshua2412-a/Diabetes_ProyectoML_202630}. Los cuatro conjuntos de datos son públicos en el repositorio UCI. La semilla principal es 42; las semillas 7 y 2026 se usan solo en el análisis de estabilidad.

\backmatter

\begin{appendices}

\section{Demostraciones}\label{app:demostraciones}

\textbf{Demostración de la Proposición~\ref{prop:cota}.} Fije $x$ y escriba $f_m=f_m(x)$, $w_m=w_m(x)$ y $\hat{f}=\sum_m w_m f_m$. Sumando y restando $\hat{f}$ dentro del cuadrado,
\begin{equation*}
\sum_{m} w_m (f_m - y)^2 = \sum_{m} w_m (f_m - \hat{f})^2 + 2(\hat{f}-y)\sum_{m} w_m (f_m - \hat{f}) + (\hat{f}-y)^2 \sum_m w_m .
\end{equation*}
Como $\sum_m w_m = 1$ y $\sum_m w_m f_m = \hat{f}$, el término central se anula y el último es $(\hat{f}-y)^2$. Reordenando se obtiene la Ecuación~\eqref{eq:ambiguedad}. Como el término de dispersión es no negativo, $\ell(\hat{f},y)\leq\sum_m w_m\,\ell(f_m,y)$, con igualdad si y solo si todas las $f_m$ con peso positivo coinciden. \hfill$\square$

\textbf{Demostración de la Proposición~\ref{prop:consistencia} (esquema).} Por la Ecuación~\eqref{eq:competencia}, $C_m(x)=\tfrac{1}{2}A^{+}(x)+\tfrac{1}{2}A^{-}(x)$, donde $A^{c}(x)$ es el promedio de $Z_i=(\tilde{f}_m(x_i)-c)^2$ sobre los $k$ vecinos de clase $c$ más cercanos a $\phi(x)$. Cada $A^{c}$ es, por tanto, un estimador de regresión por $k$ vecinos más cercanos de $E[Z\mid\phi(X)=u,\,Y=c]$ construido con los $n_c$ encuentros de esa clase. La variable $Z$ está acotada en $[0,1]$, y los pesos uniformes sobre los $k$ vecinos satisfacen las condiciones de Stone cuando $k\to\infty$ y $k/n_c\to 0$ \cite{bib_stone_1977}, lo que da consistencia en $L^1$ y, bajo la continuidad supuesta, convergencia en probabilidad para casi todo $u$. La combinación con pesos fijos $\tfrac{1}{2}$ conserva la convergencia. Dos supuestos merecen mención explícita. Primero, la independencia se exige entre pacientes y no entre encuentros, y la exclusión del propio paciente del vecindario evita que la consulta aporte información sobre sí misma. Segundo, las predicciones fuera de pliegue $\tilde{f}_m$ se tratan como fijas una vez fijada la partición de los $K$ pliegues, que es el supuesto usual del ajuste cruzado. \hfill$\square$

\end{appendices}

\bibliography{sn-bibliography}

\end{document}
